\section{Static classification and the original hypersurface}
\label{n3:static:section}

We now work entirely with the original equality data.
The preceding variational argument supplies a causal stationary field;
it does not initially supply a vacuum exterior or exclude interior points
where that field is null.
We first prove staticity where the stationary field is timelike, construct
the complete static base, and use its conformal double to exclude an
interior null set. We then reconstruct the original hypersurface,
including its boundary.
The conformal step follows the method of
Bunting and Masood-ul-Alam~\cite{BuntingMasoodUlAlam1987}.
We give the completeness and rigidity arguments needed here, without
assuming a domain-of-communications theorem or regularity of an interior
Killing horizon.

\begin{proposition}[Static classification and hypersurface reconstruction]
\label{n3:static:classification}
Let the original exterior satisfy the hypotheses of
Theorem~\ref{n3:intro:rigidity}, and let \(u,X\) be the fields supplied by
Proposition~\ref{n3:variation:adjoint}.
Then the number of boundary components is $\ell=1$, and that
component is diffeomorphic to $S^2$. Moreover,
\[
N:=u^2-|X|_g^2>0
\qquad\text{on }\operatorname{int}\Mext.
\]
The metric and function
\[
h=g+N^{-1}X^\flat\otimes X^\flat,\qquad \lambda=\sqrt N
\]
form the Schwarzschild static base of mass \(m\), with its horizon
boundary attached in the regular base coordinates described below.
There is a proper smooth spacelike embedding of the original
\(\Mext\) into maximally extended Schwarzschild spacetime of mass \(m\)
whose induced metric and future second fundamental form are \(g,K\).
Its boundary is a cross-section of the future horizon if \(u|_S>0\),
and the bifurcation sphere if \(u|_S=0\).
The embedding extends as a smooth spacelike hypersurface through \(S\),
and its chosen end approaches the corresponding spatial infinity.
\end{proposition}

We record precisely the outputs of Proposition~\ref{n3:variation:adjoint}
that will be used.
Put \(\Omega=\operatorname{int}\Mext\).
The fields \(u,X\) are smooth up to the one-sided boundary,
\(u>0\) on \(\Omega\), and \(u\geq |X|_g\).
For any fixed
\[
\frac12<q_0<\min\{q,1\},
\]
their end asymptotes have an \(O_2(r^{-q_0})\) remainder and satisfy
\begin{equation}
\label{n3:static:asymptotic-field}
(u,X)\longrightarrow(b^0,b)
 =\left(\frac Em,-\frac Pm\right),
\qquad (b^0)^2-|b|^2=1,\qquad b^0>0.
\end{equation}
On \(\R\times\Omega\), the Lorentzian metric
\begin{equation}
\label{n3:static:stationary-metric}
\mathbf g=-u^2\dd t^2+
 g_{ij}(\dd x^i+X^i\dd t)(\dd x^j+X^j\dd t)
\end{equation}
has Killing field \(\xi=\partial_t\), satisfies
\(\Ric_{\mathbf g}(\xi,\cdot)=0\), and is vacuum wherever \(N>0\).
In addition,
\begin{equation}
\label{n3:static:kid-boundary}
\sym\nabla X=-uK,\qquad
X|_S=u\nu,\qquad
\partial_\nu u+K(X,\nu)=\kappa:=\frac1{4m}.
\end{equation}
For each connected component $S_i$, either $u>0$ everywhere on $S_i$
or $u=0$ everywhere on $S_i$. The alternative may initially depend on
$i$; the same positive constant $\kappa$ occurs on all components.
These statements include the null-dust alternative before vacuum has
been proved: only the contraction \(\Ric_{\mathbf g}(\xi,\cdot)=0\)
will be used at points where \(N=0\).

\subsection{The Killing norm and the twist}
\label{n3:static:twist-subsection}

Let \(\mathcal P=\{N>0\}\subset\Omega\). On this open set define
\begin{equation}
\label{n3:static:base-definitions}
\begin{aligned}
h&=g+N^{-1}X^\flat\otimes X^\flat,
& C&=h^{-1}=g^{-1}-u^{-2}X\otimes X,\\
\lambda&=\sqrt N,
& A&=N^{-1}X^\flat,\qquad F=\dd A .
\end{aligned}
\end{equation}
Thus
\begin{equation}
\label{n3:static:threading-form}
\mathbf g=-N(\dd t-A)^2+h,\qquad
\dd V_h=\frac u{\sqrt N}\dd V_g.
\end{equation}
The tensor \(C\) extends smoothly and nonnegatively across the interior
zero set of \(N\), because \(u>0\) there.

\begin{lemma}[The original boundary has a timelike collar]
\label{n3:static:boundary-collar}
There are pairwise disjoint collars of the finitely many components
$S_1,\ldots,S_\ell$, each contained in $\mathcal P$ away from its
boundary. In the formulas for an individual collar, write $S$ for
its boundary component. If \(u|_S>0\), then
\begin{equation}
\label{n3:static:simple-norm}
\dd N|_S=2\kappa X^\flat|_S=2u\kappa\,\nu^\flat .
\end{equation}
If \(u|_S=0\), and \(s\geq0\) is the original \(g\)-normal distance
from \(S\), there are smooth \(a,Y\) such that
\begin{equation}
\label{n3:static:double-norm}
u=sa,\qquad X=s^2Y,\qquad a|_S=\kappa,\qquad
N=s^2\bigl(a^2-s^2|Y|_g^2\bigr).
\end{equation}
The AF end is also contained in \(\mathcal P\).
\end{lemma}

\begin{proof}
Fix one component and use the local notation $S$ of the statement.
On \(S\), \(X=u\nu\), hence \(N=0\) and its tangential derivatives
vanish. The normal-normal component of the KID equation gives
\[
\langle\nabla_\nu X,\nu\rangle=-uK(\nu,\nu).
\]
Consequently
\[
\partial_\nu N
=2u\partial_\nu u-2u\langle\nabla_\nu X,\nu\rangle
=2u\bigl(\partial_\nu u+uK(\nu,\nu)\bigr)
=2u\kappa.
\]
This proves Equation~\eqref{n3:static:simple-norm}, including positivity
on an exterior collar when \(u|_S>0\).

If \(u=0\) on \(S\), then \(X=0\) there and all its tangential
covariant derivatives vanish. The normal-normal and normal-tangential
components of \(\sym\nabla X=0\) imply \(\nabla_\nu X=0\) on \(S\).
Smooth division by \(s\) and then \(s^2\) yields
Equation~\eqref{n3:static:double-norm}; the boundary equation gives
\(a|_S=\kappa>0\).
Compactness of each $S_i$ makes its collar assertion uniform. Since
there are finitely many disjoint compact components, shrink their
collars to be pairwise disjoint and choose a common positive collar
width below all of these finitely many widths. Thus every original
boundary component has a timelike collar, including when the boundary
lapse alternatives differ between components.
Finally, Equation~\eqref{n3:static:asymptotic-field} gives \(N\to1\).
\end{proof}

\begin{lemma}[Norm and twist identities]
\label{n3:static:killing-identities}
Write \(\mathcal G_{\alpha\beta}=\nabla^{\mathbf g}_\alpha\xi_\beta\),
with full tensor contraction used in \(|\mathcal G|_{\mathbf g}^2\).
On \(\Omega\),
\begin{equation}
\label{n3:static:norm-wave}
\frac1u\Div_g(uC\,\dd N)=-2|\mathcal G|_{\mathbf g}^2.
\end{equation}
On \(\mathcal P\), the antisymmetric contravariant tensor
\[
Q^{ij}=uN C^{ik}C^{j\ell}F_{k\ell}
\]
satisfies
\begin{equation}
\label{n3:static:twist-divergence}
\nabla_i Q^{ij}=0.
\end{equation}
\end{lemma}

\begin{proof}
The Killing equation makes \(\mathcal G\) antisymmetric.
The differentiated Killing equation and
\(\Ric_{\mathbf g}(\xi,\cdot)=0\) imply
\(\Box_{\mathbf g}\xi=0\).
Since \(\mathbf g(\xi,\xi)=-N\), differentiation twice gives
\(\Box_{\mathbf g}N=-2|\mathcal G|_{\mathbf g}^2\).
For a \(t\)-independent scalar, the inverse metric in
Equation~\eqref{n3:static:stationary-metric} has spatial block \(C\)
and volume density \(u\sqrt{\det g}\).
Its wave operator is therefore \(u^{-1}\Div_g(uC\,\dd\cdot)\),
proving Equation~\eqref{n3:static:norm-wave}.

For completeness, put \(\theta=\dd t-A\).
The Killing one-form is \(\xi^\flat=-\lambda^2\theta\), and
\[
\xi^\flat\wedge\dd\xi^\flat=-\lambda^4\theta\wedge F.
\]
The Killing identity
\[
\dd *_\mathbf g(\xi^\flat\wedge\dd\xi^\flat)
=2*_\mathbf g\bigl(\xi^\flat\wedge\Ric_{\mathbf g}(\xi,\cdot)\bigr)
\]
has zero right side. Taking the Hodge star in the orthonormal time
form \(\lambda\theta\) shows that its left side, up to an irrelevant
overall orientation sign, is \(\dd(\lambda^3 *_hF)\).
It follows that
\[
\Div_h\bigl(\lambda^3 C^{ik}C^{j\ell}F_{k\ell}\bigr)=0.
\]
Using \(\lambda^3\dd V_h=uN\dd V_g\) gives
Equation~\eqref{n3:static:twist-divergence}.
For an antisymmetric two-tensor the extra connection term on the free
index vanishes, so this change of volume density gives exactly the
stated \(g\)-divergence.
\end{proof}

\begin{lemma}[Logarithmic control at an arbitrary null set]
\label{n3:static:log-energy}
Near the interior zero set of \(N\), put \(B=|X|_g\),
\(e=X/B\), and \(T=e^\perp\).
For every compactly supported smooth cutoff \(\eta\) in a sufficiently
small neighborhood of that set,
\begin{equation}
\label{n3:static:regularized-energy}
\sup_{\varepsilon>0}
\int u\eta^2
\frac{|\dd N|_C^2}{(N+\varepsilon)^2}\dd V_g<\infty.
\end{equation}
In particular,
\begin{equation}
\label{n3:static:transverse-log}
\int_{\mathcal P}\eta^2|\dd_T\log N|^2\dd V_g<\infty.
\end{equation}
No regularity or measure assumption on \(\{N=0\}\) is required.
\end{lemma}

\begin{proof}
On each such compact neighborhood, \(u\) and \(B\) are bounded below.
Indeed \(u>0\) in \(\Omega\) and \(B=u\) on \(\{N=0\}\).
The tensor \(C\) acts as the identity on \(T\) and has eigenvalue
\(N/u^2\) on \(e\).

We first justify a useful estimate for the smooth nonnegative function
\(N\):
\begin{equation}
\label{n3:static:smooth-nonnegative}
|\dd N|_g^2\leq C N
\end{equation}
on a compact set after slightly enlarging its neighborhood.
Choose a uniform normal-coordinate radius and a constant \(L\)
bounding the Hessian and large enough that
\(|\dd N|/L\) is smaller than that radius.
Taylor's inequality on a geodesic in direction
\(-\nabla N/|\nabla N|\), at length \(|\dd N|/L\), gives
\(0\leq N-|\dd N|^2/(2L)\).
This proves Equation~\eqref{n3:static:smooth-nonnegative}.

Let \(n=(\partial_t-X)/u\) be the future unit normal to the stationary
slices, and complete \(e\) by an orthonormal pair \(e_1,e_2\) in \(T\).
Since \(\xi=un+Be\),
\[
\mathcal G_{in}
=-\frac{N_i/2+B\mathcal G_{ie}}u.
\]
Expanding all temporal and spatial terms in the full squared norm gives
\begin{equation}
\label{n3:static:killing-square}
\begin{aligned}
-2|\mathcal G|_{\mathbf g}^2
={}&\frac{|\dd N|_g^2}{u^2}
 +\frac{4B}{u^2}\sum_{a=1}^2N_a\mathcal G_{ae}\\
&-\frac{4N}{u^2}\sum_{a=1}^2\mathcal G_{ae}^2
 -4\mathcal G_{12}^2.
\end{aligned}
\end{equation}
The fields and their derivatives are bounded on the compact
neighborhood. Equations~\eqref{n3:static:norm-wave},
\eqref{n3:static:smooth-nonnegative}, and \eqref{n3:static:killing-square}
therefore imply
\begin{equation}
\label{n3:static:norm-inequality}
\frac1u\Div_g(uC\,\dd N)
\leq C_0\bigl(N+|\dd_TN|\bigr).
\end{equation}

Multiply by \(u\eta^2/(N+\varepsilon)\) and integrate by parts.
With \(E_\varepsilon\) denoting the integral in
Equation~\eqref{n3:static:regularized-energy}, the result is
\[
E_\varepsilon
-2\int\frac{u\eta C(\dd\eta,\dd N)}{N+\varepsilon}\dd V_g
\leq C_0\int u\eta^2\dd V_g
 +C_0\int\frac{u\eta^2|\dd_TN|}{N+\varepsilon}\dd V_g .
\]
Cauchy's inequality bounds the cutoff term by
\(E_\varepsilon/4+C\int u|\dd\eta|_C^2\dd V_g\).
Since \(|\dd_TN|\leq|\dd N|_C\), the last term is bounded by
\(E_\varepsilon/4+C\int u\eta^2\dd V_g\).
This proves Equation~\eqref{n3:static:regularized-energy}.
Fatou's lemma on the positive set gives
Equation~\eqref{n3:static:transverse-log}, using the local lower bound
for \(u\).
\end{proof}

\begin{lemma}[Vanishing twist]
\label{n3:static:vanishing-twist}
The two-form \(F\) vanishes on all of \(\mathcal P\).
\end{lemma}

\begin{proof}
First note the exact contraction identity with its free index lowered
using \(g\):
\begin{equation}
\label{n3:static:twist-flux}
g_{ik}Q^{kj}A_j
=\frac1u\left(
X^j(\dd X^\flat)_{ij}
-\frac{B^2}{N}(\dd_TN)_i
\right).
\end{equation}
It follows by inserting
\[
F=N^{-1}\dd X^\flat-N^{-2}\dd N\wedge X^\flat
\quad\text{and}\quad CX=\frac N{u^2}X
\]
into \(g_{ik}Q^{kj}A_j\).
Both terms on the right of Equation~\eqref{n3:static:twist-flux}
are orthogonal to \(X\).
At \(X=0\) its second numerator is interpreted as
\(B^2\dd N-\dd N(X)X^\flat\), so no choice of \(e\) is needed there.

The end asymptotes give a constant one-form \(A_\infty\) with
\[
A-A_\infty=O_1(r^{-q_0}),\qquad
F=O(r^{-1-q_0}).
\]
Choose a smooth function equal to a linear primitive of \(A_\infty\)
near infinity and zero outside a larger end neighborhood.
Its differential \(\alpha\) is closed, equals \(A_\infty\) near
infinity, and vanishes near \(S\) and all possible interior zeros of
\(N\). These zeros are contained in a fixed compact subset of
\(\Omega\), by Lemma~\ref{n3:static:boundary-collar}.
Put \(\beta=A-\alpha\), so \(\dd\beta=F\).

For each component choose a collar cutoff $\chi_i$ that is zero near
$S_i$ and one outside its collar. Take
$\chi=\chi_{\mathrm{end}}\chi_{\mathrm{null}}
\prod_{i=1}^{\ell}\chi_i$, initially supported away from $S$,
infinity, and $\{N=0\}$. The boundary collars are pairwise disjoint;
on the support of $d\chi_i$ all other boundary factors equal one.
Testing
Equation~\eqref{n3:static:twist-divergence} by \(\chi\beta_j\) gives
\begin{equation}
\label{n3:static:twist-energy-cutoff}
\frac12\int_{\mathcal P}\chi uN
C^{ik}C^{j\ell}F_{ij}F_{k\ell}\dd V_g
=-\int_{\mathcal P}Q^{ij}(\partial_i\chi)\beta_j\dd V_g.
\end{equation}
The integrand on the left is nonnegative.

We remove the interior cutoff first, with the other two cutoffs fixed.
Let it transition from zero to one on
\(\varepsilon<N<2\varepsilon\), with derivative bounded by
\(C/\varepsilon\). Near that band \(\beta=A\).
Equation~\eqref{n3:static:twist-flux}, its orthogonality to \(X\),
and \(N\asymp\varepsilon\) bound its contribution by
\[
C\int_{\{\varepsilon<N<2\varepsilon\}}
\bigl(1+|\dd_T\log N|^2\bigr)\dd V_g
\]
on a fixed compact set.
This tends to zero by Lemma~\ref{n3:static:log-energy}.
Indeed the bands tend pointwise to the empty set within \(\mathcal P\),
and their indicators are dominated by an integrable function there.
This reasoning remains valid when \(\{N=0\}\) has positive measure.

Next remove the boundary cutoffs, keeping the end cutoff fixed.
In the collar of $S_i$ let $s$ be the original normal distance and
let $\chi_i$ transition on $s\in(\delta,2\delta)$ with derivative
bounded by $C_i/\delta$. For this local computation write $S=S_i$.
If \(u|_S>0\), then
\(N\asymp s\), \(X=u\nu+O(s)\), and \(\dd_TN=O(s)\).
The flux in Equation~\eqref{n3:static:twist-flux} is bounded.
Its orthogonality to \(X\) makes its \(g\)-normal component \(O(s)\),
so multiplication by the cutoff derivative \(O(\delta^{-1})\)
and integration over the collar gives a term tending to zero.
If \(u|_S=0\), Equation~\eqref{n3:static:double-norm} gives
\[
u\asymp s,\quad X=O(s^2),\quad \dd X^\flat=O(s),
\quad N\asymp s^2,\quad \dd N=O(s).
\]
The entire flux is then \(O(s^2)\), with the same conclusion.
Each boundary error therefore tends to zero independently of which
lapse alternative holds on that component. The derivative of the
finite product is the finite sum of these collar contributions;
letting $\delta\downarrow0$ makes their sum tend to zero.

Finally use a radial cutoff on \(R<r<2R\).
Here \(Q=O(r^{-1-q_0})\) and \(\beta=O(r^{-q_0})\);
its contribution is \(O(R^{1-2q_0})\), which tends to zero.
Fatou's lemma in Equation~\eqref{n3:static:twist-energy-cutoff}
now makes its nonnegative integrand vanish everywhere on
\(\mathcal P\).
Since \(uN>0\) and \(C\) is positive definite there, \(F=0\).
\end{proof}

\subsection{The complete static base}
\label{n3:static:base-subsection}

Let $\mathcal E$ be the component of $\mathcal P$ containing the
connected far AF end. We first derive the static equations on every
component of $\mathcal P$, then prove that there is only one component.
Since \(F=0\), the form \(A\) is locally exact.
Equation~\eqref{n3:static:threading-form} can locally be written as
\(-\lambda^2\dd T^2+h\).
The Christoffel symbols involving \(T\) are
\(\Gamma^T_{Ti}=\lambda_i/\lambda\) and
\(\Gamma^i_{TT}=\lambda h^{ij}\lambda_j\).
The vacuum equations consequently give, on all of $\mathcal P$,
\begin{equation}
\label{n3:static:static-equations}
\lambda\Ric_h=\Hess_h\lambda,\qquad
\Delta_h\lambda=0,\qquad \Scal_h=0.
\end{equation}
These equations are global even though no global primitive of \(A\)
has yet been constructed.

\begin{lemma}[Connected positive set]
\label{n3:static:connected-positive-set}
The set $\mathcal P$ equals its infinity component $\mathcal E$.
In particular, every original boundary collar belongs to $\mathcal E$.
\end{lemma}

\begin{proof}
The unique end of $\Mext$ lies in $\mathcal P$ outside a compact set.
Any component $V$ of $\mathcal P$ other than $\mathcal E$ therefore
has compact closure in the original closed exterior $\Mext$.
Its frontier there lies in $\{N=0\}\cup S$: an interior frontier
point with $N>0$ has a connected positive neighborhood meeting $V$
and hence belongs to $V$, a contradiction. The continuous function
$\lambda=\sqrt N$ is zero on that frontier. The frontier is nonempty;
otherwise $V$ would be open and closed in the connected open exterior
$\Omega$, which also contains the far end in $\mathcal E$.

Choose any point of $V$, where $\lambda>0$. Compactness of the
closure and the zero frontier values imply that the positive maximum
of $\lambda$ is attained at an interior point of $V$. The strong
maximum principle for the smooth locally elliptic equation
$\Delta_h\lambda=0$ makes $\lambda$ constant on $V$, contradicting
its zero frontier values. Thus no such component exists.
Every boundary collar from Lemma~\ref{n3:static:boundary-collar} is a
subset of $\mathcal P=\mathcal E$, as claimed.
\end{proof}

\begin{lemma}[Completion and the lapse range]
\label{n3:static:base-completion}
The function \(\lambda\) satisfies \(0<\lambda<1\) on \(\mathcal E\).
Approach to an interior zero of \(N\) has infinite \(h\)-length.
Attaching all components of $S$ gives a regular base boundary on which
\begin{equation}
\label{n3:static:base-boundary-data}
h|_{TS}=g|_{TS},\qquad
\lambda=0,\qquad
\partial_{\nu_h}\lambda=\kappa,\qquad
\mathrm{II}_h=0.
\end{equation}
At a component $S_i$ with $u|_{S_i}>0$, the base metric admits a
smooth reflection with odd lapse. At a component with $u|_{S_i}=0$,
it admits a $C^2$ reflection with odd lapse. These componentwise
attachments form a complete base with compact boundary, and its
two-copy double is a connected orientable complete manifold.
\end{lemma}

\begin{proof}
Near an interior zero set, \(B,u\) have positive lower bounds.
Contracting \(F=0\) with \(X\) in
Equation~\eqref{n3:static:twist-flux} gives
\[
|\dd_T\log N|\leq B^{-2}|X|\,|\dd X^\flat|\leq C.
\]
Together with Equation~\eqref{n3:static:smooth-nonnegative}, this yields
\begin{equation}
\label{n3:static:complete-log-bound}
|\dd\log N|_h^2
=|\dd_T\log N|^2+\frac{|\dd_eN|^2}{u^2N}\leq C
\end{equation}
on compact neighborhoods of the interior zero set.
A finite-length path approaching that set would make
\(\log N\to-\infty\), contradicting this bound.

Every possible limit of \(\mathcal E\) in the compact part of the
original closed exterior either belongs to \(\mathcal E\), lies in
\(\{N=0\}\), or belongs to \(S\). At the latter two types of points
\(\lambda\to0\), while \(\lambda\to1\) at the AF end.
If \(\lambda>1\) somewhere, a positive superlevel set bounded away
from one is contained in a compact subset of \(\mathcal E\).
The harmonic function would attain an interior maximum, forcing it
constant, contrary to its limit at infinity.
Thus \(\lambda\leq1\).
If equality holds at one interior point, the strong maximum principle
gives \(\lambda\equiv1\) on \(\mathcal E\).
That component then has no interior boundary, since \(N\) remains one
at any such boundary point. It is open and closed in connected
\(\Omega\), hence equals \(\Omega\), contradicting the timelike collar
where \(\lambda\to0\) at \(S\).
Therefore \(0<\lambda<1\).

For boundary regularity fix a component $S_i$ and write $S$ for it
through the following local calculation. Suppose first that \(u|_S>0\).
By Equation~\eqref{n3:static:simple-norm}, use \((N,y^1,y^2)\)
as original smooth collar coordinates. Smooth division gives
\[
X^\flat=a(N,y)\dd N+N\beta_A(N,y)\dd y^A,\qquad
a(0,y)=\frac1{2\kappa}.
\]
Writing \(N=\lambda^2\), the coefficients of \(h\) are
\begin{equation}
\label{n3:static:even-base}
\begin{aligned}
h_{\lambda\lambda}&=4\lambda^2g_{NN}+4a^2,\\
h_{\lambda A}&=2\lambda(g_{NA}+a\beta_A),\\
h_{AB}&=g_{AB}+\lambda^2\beta_A\beta_B,
\end{aligned}
\end{equation}
with all right-side coefficients evaluated at \((\lambda^2,y)\).
They extend smoothly to negative \(\lambda\).
The map \((\lambda,y)\mapsto(-\lambda,y)\) is an isometry:
the diagonal blocks are even and the mixed block is odd.
At zero, \(h_{\lambda\lambda}=\kappa^{-2}\),
\(h_{\lambda A}=0\), and \(h_{AB}=g_{AB}|_S\).
The fixed hypersurface of this reflection is totally geodesic,
and all assertions in Equation~\eqref{n3:static:base-boundary-data}
follow.

If \(u|_S=0\), put \(D=a^2-s^2|Y|^2\) in
Equation~\eqref{n3:static:double-norm}. Then
\[
h=g+s^2D^{-1}Y^\flat\otimes Y^\flat,\qquad
\lambda=s\sqrt D,\qquad D|_S=\kappa^2.
\]
These are smooth in the original one-sided collar.
The static equations extend by continuity to give
\(\Hess_h\lambda=0\) at \(S\).
Its tangential components are
\(\kappa\,\mathrm{II}_h\), so \(\mathrm{II}_h=0\).
In \(h\)-Gaussian coordinates \(\rho\geq0\), write
\(h=\dd\rho^2+h_\rho\). Then
\(\partial_\rho h_\rho|_0=0\), and
\(\partial_\rho^2\lambda|_0=0\).
Even reflection of \(h_\rho\) and odd reflection of \(\lambda\)
are therefore \(C^2\).

These constructions apply on the finitely many disjoint collars,
with no compatibility condition between different lapse alternatives.
They attach the entire original $S$, by
Lemma~\ref{n3:static:connected-positive-set}. Let $\overline{\mathcal E}$
denote this attached base and let $j:\overline{\mathcal E}\to\Mext$
be its identity map on the interior and on boundary points. It is a
homeomorphism onto the original exterior with any interior zero set
removed, but need not be a boundary diffeomorphism: on a positive-lapse
collar the original coordinate is $N=\lambda^2$. The identity
$h=g+N^{-1}X^\flat\otimes X^\flat$ gives
$h\geq j^*g$ in the interior and, by continuity in each regular collar,
on the attached boundary as a nonnegative quadratic-form inequality.

It remains to check completeness without classifying the missing
sets. The original exterior is complete as a metric space with
boundary: a finite-length Cauchy path is Cauchy in the complete
ambient initial-data manifold and its limit remains in the closed
exterior. The inequality $h\geq j^*g$ makes the image of a finite
$h$-length escaping path have such an original limit.
An interior limit with \(N>0\) is an ordinary smooth base limit.
An interior zero is impossible by
Equation~\eqref{n3:static:complete-log-bound}.
A limit in \(S\) is accounted for by the regular collar just
constructed. The AF end has infinite distance.
These possibilities exhaust finite-length escapes, proving completeness
of $\overline{\mathcal E}$ with its boundary included.

Take two copies of this base and identify corresponding points of
each $S_i$ by the collar reflections just constructed. The resulting
double is connected because the base is connected and its boundary
is nonempty. Give the second copy the opposite orientation; the
identifications yield an orientable manifold without boundary.
There is no assertion that this double is simply connected.
The reflected metric $h_{\mathrm d}$ and odd reflected lapse are
$C^2$ across every join, and smooth at the positive-lapse joins.
The static scalar and harmonic equations hold continuously across
the joins, since the indicated second jets match.

For completeness, fold the double onto $\overline{\mathcal E}$ by
identifying the two copies. The fold is distance nonincreasing,
because lengths on either side are the base lengths. A Cauchy
sequence in the double projects to a Cauchy sequence in the complete
base. The base is a complete locally compact length space, hence is
proper, so the projected sequence lies in a compact base set. The
preimage of that set is a quotient of two compact copies and is
compact. The original sequence consequently has a convergent
subsequence and, being Cauchy, converges. This proves completeness
of the double.
\end{proof}

\subsection{Static asymptotics and the compactified end}
\label{n3:static:asymptotic-subsection}

After a constant linear change of the original end coordinates,
\(h=\delta+O_2(r^{-q_0})\) and
\(\lambda=1+O_2(r^{-q_0})\).
Indeed its original constant metric is
\(I+b\otimes b\), by Equation~\eqref{n3:static:asymptotic-field}.
We prove the improvement required for conformal doubling explicitly;
the harmonic-coordinate mechanism agrees with
Bartnik~\cite[Theorem 3.1, Proposition 3.3, and Theorem 4.3]{Bartnik1986}.

\begin{lemma}[Harmonic-coordinate expansion]
\label{n3:static:asymptotics}
There are asymptotically Cartesian coordinates on the end and a
number \(0<\epsilon<1\) such that, with
\(\gamma=\lambda^2h\) and \(U=\log\lambda\),
\begin{equation}
\label{n3:static:gamma-expansion}
\gamma_{ij}=\delta_{ij}+O_k(r^{-1-\epsilon}),\qquad
U=-\frac Mr+\frac{b_1\cdot x}{r^3}
       +O_k(r^{-2-\epsilon})
\end{equation}
for every fixed finite \(k\).
Here \(M>0\); the vector \(b_1\) is unrelated to the shift asymptote
\(b\).
Only the original two metric derivatives and the stated adjoint
decay are needed to start this improvement.
\end{lemma}

\begin{proof}
The three-dimensional conformal formulas applied to
Equation~\eqref{n3:static:static-equations} give
\begin{equation}
\label{n3:static:gamma-equations}
\Ric_\gamma=2\dd U\otimes\dd U,\qquad \Delta_\gamma U=0.
\end{equation}
Initially \(\gamma-\delta,U=O_2(r^{-q_0})\).
We give the coordinate construction with this derivative count.

Move to a sufficiently distant end and extend \(\gamma\) to a
metric on \(\R^3\), equal to \(\delta\) inside a large ball,
using a fixed-ratio cutoff annulus.
In the resulting coordinates write
\[
\Delta_\gamma=(\delta^{ab}+a^{ab})\partial_{ab}+c^a\partial_a .
\]
The coefficients satisfy \(a=O_2(r^{-q_0})\),
\(c=O_1(r^{-1-q_0})\), and their scaled Hölder norms through
the orders needed here are small after increasing the cutoff radius.
The bounded second derivatives in the hypothesis give scaled
Lipschitz bounds on first derivatives, so they suffice for any fixed
Hölder exponent less than one at this stage.

Put \(s_0=1-q_0\in(0,1/2)\).
For a source with weighted scaled \(C^{0,\alpha}\) size
\(O(r^{s_0-2})\), the Euclidean inverse normalized to vanish at zero is
\begin{equation}
\label{n3:static:subtracted-newton}
Pf(x)=-\frac1{4\pi}\int_{\R^3}
\left(\frac1{|x-y|}-\frac1{|y|}\right)f(y)\dd y.
\end{equation}
It maps that weighted norm boundedly to scaled
\(C^{2,\alpha}\) size \(O(r^{s_0})\).
To verify the zeroth-order bound, split at \(|y|=2|x|\).
On the far part, the kernel difference is bounded by
\(C|x||y|^{-2}\), and its integral is \(O(|x|^{s_0})\)
because \(s_0<1\).
On the near part the separate kernels have the same bound because
\(s_0>0\).
Interior Poisson estimates on rescaled annuli give the derivative
and Hölder bounds.
The ordinary local estimates also apply on the fixed inner ball.

The equations
\[
v^i=P\bigl(-c^i-a^{ab}\partial_{ab}v^i-c^a\partial_av^i\bigr)
\]
are contractions in that weighted space.
Indeed the operator norm of the last two terms is bounded by a
constant times the small scaled norms of \(a\) and \(rc\).
They produce harmonic coordinates \(x^i+v^i\) with
\(v=O_{C^{2,\alpha}}(r^{1-q_0})\) and
\(Dv=O(r^{-q_0})\).
These functions are coordinates sufficiently far out.

There is a derivative issue in changing metric coordinates which
must be addressed before using a classical Ricci equation.
The original assumptions bound \(a,c\), respectively, in scaled
\(W^{2,\infty}\) and \(W^{1,\infty}\).
Differentiate the equation for \(v\) once and apply local
\(W^{2,p}\) estimates to \(Dv\) on rescaled annuli.
All differentiated coefficients and forcing are controlled by
these bounds and the preceding \(C^{2,\alpha}\) estimate.
For every finite \(p\) this gives
\[
v=O_{W^{3,p}}(r^{1-q_0}).
\]
Consequently the transformed \(\gamma-\delta,U\) have scaled
\(W^{2,p}\) size \(O(r^{-q_0})\).
Take \(p>3\); Morrey embedding gives scaled \(C^{1,\alpha}\)
control for some \(\alpha>0\).
In the harmonic chart Equation~\eqref{n3:static:gamma-equations}
has the elliptic form
\[
-\tfrac12\gamma^{ab}\partial_{ab}\gamma_{ij}
 +Q_{ij}(\gamma^{-1},D\gamma)=2U_iU_j,\qquad
\gamma^{ab}\partial_{ab}U=0,
\]
where \(Q_{ij}\) is quadratic in \(D\gamma\).
Schauder estimates first give scaled \(C^{2,\alpha}\) control;
differentiating this elliptic system then gives
\(\gamma-\delta,U=O_k(r^{-q_0})\) for each finite \(k\).
No higher original-coordinate derivative bound has been assumed.

We use two elementary exterior Poisson estimates.
If \(w=o(1)\), \(\Delta_\delta w=f\), and
\(f=O_k(r^{-2-t})\), \(1<t<2\), then
\begin{equation}
\label{n3:static:poisson-monopole}
w=\frac ar+O_k(r^{-t}).
\end{equation}
To prove this, cut \(w\) off on a fixed inner annulus and extend it
by zero; this changes \(f\) only on a compact set.
The extended solution equals the Newton potential of its Laplacian,
since their difference is an entire decaying harmonic function.
The coefficient is \(a=-(4\pi)^{-1}\int f\).
For \(|y|<r/2\), subtracting \(1/r\) from the kernel costs
\[
Cr^{-2}\int_{|y|<r/2}|y||f(y)|\dd y=O(r^{-t}).
\]
The remaining region, including the integrable kernel singularity,
has the same bound. Rescaled interior estimates prove the derivative
version.
If instead \(f=O_k(r^{-4-\epsilon})\), \(0<\epsilon<1\), its first
moments converge and the same reasoning gives
\begin{equation}
\label{n3:static:poisson-dipole}
w=\frac ar+\frac{d\cdot x}{r^3}+O_k(r^{-2-\epsilon}).
\end{equation}
Here subtract \(1/r+(x\cdot y)/r^3\) in the near region.
The remainder costs
\(Cr^{-3}\int_{|y|<r/2}|y|^2|f(y)|\dd y
 =O(r^{-2-\epsilon})\); the far region has that order as well.

Set \(\epsilon=2q_0-1\in(0,1)\).
The harmonic-coordinate equations and initial symbol bounds yield
\[
\Delta_\delta(\gamma_{ij}-\delta_{ij}),
\ \Delta_\delta U=O_k(r^{-2-2q_0}).
\]
Equation~\eqref{n3:static:poisson-monopole} gives
\[
\gamma_{ij}=\delta_{ij}+\frac{A_{ij}}r
 +O_k(r^{-1-\epsilon}),\qquad
U=\frac ar+O_k(r^{-1-\epsilon}).
\]
The harmonic-coordinate condition is
\(\partial_j(\sqrt{\det\gamma}\,\gamma^{ij})=0\).
Expanding it at the displayed order gives
\[
\left(A_{ij}-\tfrac12(\tr A)\delta_{ij}\right)
\frac{x^j}{r^3}=O(r^{-2-\epsilon}).
\]
Taking a radial limit shows \(A=\tfrac12(\tr A)I\).
Its trace in dimension three is zero, hence \(A=0\).
Since \(D^2U=O_k(r^{-3})\), the scalar equation now improves to
\[
\Delta_\delta U
=(\delta^{ab}-\gamma^{ab})\partial_{ab}U
=O_k(r^{-4-\epsilon}).
\]
Equation~\eqref{n3:static:poisson-dipole} proves
Equation~\eqref{n3:static:gamma-expansion}, initially with a real \(M\).

Its positivity uses the lapse range from
Lemma~\ref{n3:static:base-completion}.
The function \(-U>0\) is \(\gamma\)-harmonic.
For \(0<\eta<1\) and a fixed positive \(C\), the function
\[
w=r^{-1}+Cr^{-1-\eta}
\]
is positive and strictly \(\gamma\)-subharmonic sufficiently far out:
\(\Delta_\gamma r^{-1}=O(r^{-4-\epsilon})\), while the leading
Laplacian of \(r^{-1-\eta}\) is
\(\eta(1+\eta)r^{-3-\eta}\).
Choose \(c>0\) so small that \(-U\geq cw\) on a large inner sphere.
The difference tends to zero and is superharmonic, so the exterior
minimum principle gives \(-U\geq cw\) throughout that end.
Taking the radial limit proves \(M\geq c>0\).
\end{proof}

\begin{lemma}[The two conformal metrics]
\label{n3:static:conformal-completion}
Define
\[
h_\pm=\left(\frac{1\pm\lambda}{2}\right)^4h
\quad\text{on }\mathcal E.
\]
Both are scalar flat. The plus end satisfies
\(h_+=\delta+O_k(r^{-1-\epsilon})\), hence has zero ADM mass.
The minus end compactifies by one point to a nondegenerate
\(C^{1,\epsilon}\cap W^{2,p}_{\mathrm{loc}}\) metric for some
\(p>3\), with scalar curvature zero as a distribution at that point.
\end{lemma}

\begin{proof}
Scalar flatness follows from
Equation~\eqref{n3:static:static-equations} and the three-dimensional
conformal scalar-curvature formula, since \(1\pm\lambda\) are harmonic.
In terms of the preceding variables,
\begin{equation}
\label{n3:static:conformal-gamma}
h_+=\cosh^4(U/2)\gamma,\qquad
h_-=\sinh^4(U/2)\gamma.
\end{equation}
The first factor is \(1+O_k(r^{-2})\), proving the assertion for
\(h_+\). Its ADM flux is \(O(r^{-\epsilon})\), so its mass is zero.

For the minus end invert \(x=y/|y|^2\), write \(\rho=|y|\), and
put \(R(y)=I-2y\otimes y/\rho^2\).
The inversion differential is \(\rho^{-2}R(y)\), with
\(R(y)^TR(y)=I\).
Equation~\eqref{n3:static:gamma-expansion} gives
\[
I^*\gamma=\rho^{-4}
\bigl(I+O_k(\rho^{1+\epsilon})\bigr),
\qquad
U(I(y))=-M\rho+(b_1\cdot y)\rho
 +O_k(\rho^{2+\epsilon}).
\]
Derivatives of the angular matrix \(R(y)\) cost the corresponding
powers of \(\rho^{-1}\), so the displayed metric estimate retains
the symbol bounds. Dividing the expansion of \(\sinh(U/2)\) by
\(\rho\), and using \(\epsilon<1\), gives
\begin{equation}
\label{n3:static:inverted-minus}
I^*h_-=\left(\frac M2\right)^4
\left[
\left(1-\frac{4b_1\cdot y}{M}\right)I
 +O_k(\rho^{1+\epsilon})
\right].
\end{equation}
In particular, its only first-order term is linear and scalar.

For the remainder \(E(y)\), the bounds
\[
|E|\leq C\rho^{1+\epsilon},\quad
|DE|\leq C\rho^\epsilon,\quad
|D^2E|\leq C\rho^{\epsilon-1}
\]
show that \(E(0)=DE(0)=0\) extends it as \(C^{1,\epsilon}\).
To see the Hölder estimate between two points, use the gradient
bound when their separation is comparable to their radii; otherwise
integrate the Hessian along the short segment in a common annulus.
Moreover choose
\[
3<p<\frac3{1-\epsilon}.
\]
Then \(D^2E\in L^p\). Integration by parts across small coordinate
spheres has an error bounded by a constant times their area,
because the first derivatives are bounded.
It follows that the punctured classical derivatives are the weak
derivatives of the extension.
The Christoffel symbols are continuous with weak derivatives in
\(L^p\), so scalar curvature is its usual \(L^p\) expression
in \(\partial\Gamma\) and \(\Gamma\Gamma\).
It vanishes off the point and therefore also as a distribution on
the whole ball. There is no point curvature term.
\end{proof}

\subsection{Zero-mass rigidity with arbitrary remaining ends}
\label{n3:static:spin-subsection}

The completed conformal manifold can still have ends corresponding
to an interior zero set. We need a rigidity statement that permits
these ends, and also the regularity in
Lemma~\ref{n3:static:conformal-completion}.
The smooth arbitrary-end positive mass theorem is established in
\cite[Theorem B]{CecchiniZeidler2024}.
In our scalar-flat situation the following direct spinor proof also
handles the compactified point. It uses Witten's
identity~\cite{Witten1981}; no asymptotic condition is imposed on any
other end.

\begin{lemma}[Scalar-flat rigidity at a faster-decaying end]
\label{n3:static:spin-rigidity}
Let \((W,\mathfrak g)\) be a connected orientable complete
three-dimensional Riemannian manifold without boundary.
Assume that its metric is locally \(C^{1,\alpha}\cap W^{2,p}\),
for some \(\alpha>0,p>3\), is scalar flat as a distribution,
and has an end on which it is smooth with
\[
\mathfrak g=\delta+O_2(r^{-1-\epsilon}),\qquad \epsilon>0.
\]
The other ends can be arbitrary.
Then \((W,\mathfrak g)\) is isometric to Euclidean space.
The asserted regularity is sufficient at a compactified point and
at a reflected hypersurface.
\end{lemma}

\begin{proof}
An orientable three-manifold is spin.
Choose a spin structure and use a smooth background metric to
identify its spinor bundle with that for \(\mathfrak g\).
The positive square-root change of orthonormal frames has the
stated metric regularity.
The spin connection consequently has continuous locally bounded
coefficients, with weak first derivatives in \(L^p_{\mathrm{loc}}\).
Write \(\nabla^S\) for that connection, \(D\) for the Dirac operator,
and \(c\) for Clifford multiplication.

The weak Lichnerowicz formula has the following integrated form for
every compactly supported \(H^1\) spinor \(\phi\):
\begin{equation}
\label{n3:static:lichnerowicz}
\|D\phi\|_2^2=\|\nabla^S\phi\|_2^2.
\end{equation}
Here and below norms and integrals are for \(\mathfrak g\).
To check the formula at the indicated regularity, first use a smooth
compactly supported spinor. Approximate the metric on its support
in \(C^1\) and \(W^{2,p}\) by smooth metrics, using the same
background-bundle identification.
The smooth Lichnerowicz formula passes to the limit: connection
coefficients converge uniformly and scalar curvatures converge
in \(L^p\).
The limiting scalar term is zero by the assumed distributional
scalar flatness and the \(L^p\) curvature expression.
Density then proves Equation~\eqref{n3:static:lichnerowicz} for compact
\(H^1\) spinors.
Thus no integration surface is retained around the compactified point.

We next establish the coercivity needed to complete compact spinors.
For a compactly supported scalar function \(f\) on the AE end,
one-dimensional integration by parts along coordinate rays gives
\begin{equation}
\label{n3:static:radial-hardy}
\int_{r_0}^\infty f^2\dd r
\leq4\int_{r_0}^\infty r^2|\partial_rf|^2\dd r.
\end{equation}
The inner boundary term has the favorable sign.
Apply this to \(f=|\phi|\) and integrate over angles.
Metric comparability and
\(|\dd|\phi||\leq|\nabla^S\phi|\) give
\begin{equation}
\label{n3:static:spin-hardy}
\int_{\mathrm{AE}}r^{-2}|\phi|^2\dd V_{\mathfrak g}
\leq C\|\nabla^S\phi\|_2^2.
\end{equation}
For any fixed compact \(K\subset W\), connect it to a fixed end
annulus by a relatively compact connected domain.
Poincaré's inequality with the norm on that annulus retained gives
\[
\||\phi|\|_{L^2(K)}
\leq C_K\left(
\|\dd|\phi|\|_{L^2(\text{domain})}
 +\|\phi\|_{L^2(\text{annulus})}\right).
\]
For example, this version follows from the ordinary Poincaré
inequality after subtracting the mean; the observed annulus
controls that mean because it has positive measure.
Equation~\eqref{n3:static:spin-hardy} controls its last term.
We have proved
\begin{equation}
\label{n3:static:local-coercivity}
\|\phi\|_{L^2(K)}\leq C_K\|\nabla^S\phi\|_2.
\end{equation}
Connectedness is the only global input in this propagation.

Let \(\mathcal H\) be the completion of compactly supported smooth
spinors in the norm \(\|\nabla^S\phi\|_2\).
Equation~\eqref{n3:static:local-coercivity} embeds \(\mathcal H\)
in \(L^2_{\mathrm{loc}}\), and the locally bounded connection
then embeds it in \(H^1_{\mathrm{loc}}\).
Both \(\nabla^S:\mathcal H\to L^2\) and
\(D:\mathcal H\to L^2\) are well-defined, and
Equation~\eqref{n3:static:lichnerowicz} makes the latter an isometry.
Its range is therefore closed.
Equation~\eqref{n3:static:spin-hardy} also passes to this completion.

Choose an arbitrary constant spinor \(z\) in the ordinary
orthonormalized coordinate frame on the distinguished end.
Let \(\psi_0\) equal \(z\) near infinity and vanish before a
slightly smaller end neighborhood.
Because the connection there is \(O(r^{-2-\epsilon})\),
\(\nabla^S\psi_0,D\psi_0\in L^2\).
Orthogonally project \(D\psi_0\) onto the closed subspace
\(D\mathcal H\) and choose \(\eta\in\mathcal H\) such that
\[
D\eta=-\operatorname{proj}_{D\mathcal H}(D\psi_0).
\]
Set \(\psi=\psi_0+\eta\), \(w=D\psi\).
Then \(w\in L^2\) and
\[
\int\langle w,D\phi\rangle\dd V_{\mathfrak g}=0
\quad\text{for every compact smooth }\phi.
\]
Thus \(Dw=0\) weakly.
Local elliptic regularity for a first-order elliptic operator with
locally Lipschitz principal coefficients and bounded lower-order
coefficients gives \(w\in H^1_{\mathrm{loc}}\).
This regularity can also be seen by freezing the principal
coefficients on small coordinate balls: the constant-coefficient
Dirac estimate controls one derivative in \(L^2\), the oscillation
is absorbed, and commutators of mollification are bounded by the
local Lipschitz norm. Apply the estimate to regularized \(w\)
and pass to a weak limit.
The metric in the lemma has more than this coefficient regularity.

Completeness supplies compactly supported Lipschitz distance cutoffs
\(\chi_R\), equal to one on the radius-\(R\) ball and with
\(|\dd\chi_R|\leq C/R\).
Indeed a complete locally compact Riemannian length space is proper,
and a piecewise linear cutoff of its distance has these properties.
Equation~\eqref{n3:static:lichnerowicz} applies to \(\chi_Rw\), and
the weak equation gives
\[
\|\nabla^S(\chi_Rw)\|_2^2
=\|D(\chi_Rw)\|_2^2
=\|c(\dd\chi_R)w\|_2^2
\leq \frac C{R^2}\|w\|_2^2.
\]
On every fixed compact set the cutoff eventually equals one.
It follows that \(\nabla^S w=0\).
A parallel spinor has constant length on the connected manifold;
the distinguished end has infinite volume, so \(w\in L^2\)
forces \(w=0\).
No condition at another end entered this argument.

It remains to show that \(\psi\), rather than just \(D\psi\), is
parallel. Define the continuous sesquilinear form
\[
\mathcal B(a,b)=
\int\left(\langle\nabla^Sa,\nabla^Sb\rangle
-\langle Da,Db\rangle\right)\dd V_{\mathfrak g}
\]
on fields for which the indicated derivatives are in \(L^2\).
For compact \(\phi\), the polarized Lichnerowicz identity,
with one slot compactly supported, gives
\(\mathcal B(\psi_0,\phi)=0\).
Approximate \(\eta\) in \(\mathcal H\) by compact spinors.
Continuity and Equation~\eqref{n3:static:lichnerowicz} imply
\[
\mathcal B(\psi_0,\eta)=0,\qquad
\mathcal B(\eta,\eta)=0,\qquad
\mathcal B(\psi,\psi)=\mathcal B(\psi_0,\psi_0).
\]
The last expression is zero. To see this directly, integrate the
Lichnerowicz identity for \(\psi_0\) up to a large end sphere.
There is no contribution at another end because \(\psi_0\)
vanishes there. On that sphere \(\psi_0=O(1)\) and
\(\nabla^S\psi_0,D\psi_0=O(r^{-2-\epsilon})\);
the boundary integral is \(O(r^{-\epsilon})\) and tends to zero.
This is also the zero ADM boundary term of the spin proof.
Since \(D\psi=w=0\), we conclude that
\(\|\nabla^S\psi\|_2^2=0\).

For \(z\ne0\) the resulting parallel spinor cannot be zero.
Otherwise \(\eta=-\psi_0\), contradicting
Equation~\eqref{n3:static:spin-hardy}, because a nonzero constant
spinor has infinite \(r^{-2}\)-weighted \(L^2\) norm on the end.
The construction is linear in \(z\), so it produces a full
complex basis of parallel spinors at each point.
Their parallel equations first imply \(C^{1,\beta}\) regularity
locally for some \(\beta>0\): from \(H^1\subset L^6\) and the
bounded connection one obtains \(W^{1,6}\), hence local Hölder
continuity. The connection coefficients are locally Hölder by the
metric regularity, so the parallel equations make the first
derivatives locally Hölder as well.
Their curvature equations hold weakly.
The spin representation is faithful on the Lie algebra
\(\mathfrak{spin}(3)\), so a full parallel spinor frame makes
the Riemann curvature zero.
Equivalently, its spinor bilinears give a global parallel
orthonormal tangent frame.
The dual coframe is closed by torsion-freeness and integrates
locally to coordinates in which the metric is exactly Euclidean.
This also removes any apparent metric singularity at the
compactified point.

Finally, completeness makes the universal flat cover Euclidean
\(\R^3\). The global parallel frame makes its deck transformations
translations. A nontrivial discrete translation group has rank
at least one, and a quotient by it has volume growth at most
quadratic. The distinguished AE end gives a cubic lower bound:
radial paths in that end put a Euclidean annulus of radius
comparable to \(R\) inside a metric ball of radius \(CR\).
The two volume bounds are incompatible.
The deck group is therefore trivial, proving the lemma.
\end{proof}

\begin{lemma}[The infinity component has no interior missing boundary]
\label{n3:static:no-internal-null}
The component \(\mathcal E\) is all of \(\Omega\).
In particular \(N>0\) everywhere in the open exterior and the
stationary metric is vacuum there.
\end{lemma}

\begin{proof}
The complete base of Lemma~\ref{n3:static:base-completion} reaches every
component of $S$. Take its complete two-copy double, use $h_+$ on
the plus copy and $h_-$ on the minus copy, and identify corresponding
boundary points. The signed reflected lapse expresses the joined
metric as
\[
q=\left(\frac{1+\lambda_{\mathrm{signed}}}{2}\right)^4h_{\mathrm d}.
\]
It is smooth at each positive-lapse join and $C^2$ at each zero-lapse
join. The matched second jets make its scalar curvature zero across
every join, as well as on the two open sides. Compactify the unique
AF end of the minus copy by
Lemma~\ref{n3:static:conformal-completion}.

We verify completeness of this conformal manifold without placing any
restriction on the number or regularity of the possible internal-null
ends. On the plus copy the factor $(1+\lambda)/2$ is at least $1/2$.
Fix a sufficiently large original AF tail. On its complement the
continuous function $\sqrt N$ is defined on a compact subset of the
original closed exterior, equals zero on $S$ and the internal zero
set, and is strictly less than one elsewhere by
Lemma~\ref{n3:static:base-completion}. Its maximum on this compact set
is therefore some $a_0<1$. Thus on the entire corresponding minus
remainder $(1-\lambda)/2\geq(1-a_0)/2>0$ uniformly. On the double
away from the minus AF tail the conformal metric is consequently
bounded below by a fixed positive multiple of the complete doubled
base metric. All joins are already attached. A finite-length escape
can therefore only run down the minus AF tail, and its proved
nondegenerate one-point compactification supplies its limit. More
explicitly, truncate that tail at a large sphere. The other side of
the sphere is complete with compact boundary by the uniform lower
bound, and the compactified tail is a compact metric region with
that same boundary; their gluing is complete.

The conformal manifold is connected, orientable, and without boundary.
It has local $C^{1,\alpha}\cap W^{2,p}$ regularity for some
$\alpha>0$ and $p>3$: the finitely many joins are at least $C^2$,
and the sole compactification point has the regularity established
in Lemma~\ref{n3:static:conformal-completion}. Its scalar curvature is
zero distributionally, including at that point, and its distinguished
plus end has the fast decay required by
Lemma~\ref{n3:static:spin-rigidity}. That lemma makes this entire
manifold Euclidean space.

Choose a sufficiently large sphere in the distinguished plus end.
Under the Euclidean isometry it is a compact embedded sphere. Its
end tail is the unbounded component of the complement: paths from
large end radii to that sphere have lengths tending to infinity,
and the tail has precisely that sphere as frontier. The other
component is bounded and has compact closure. A sequence in
$\mathcal E$ approaching an interior point with $N=0$ lies outside
the chosen end tail and escapes every compact set of the complete
conformal manifold. Indeed it cannot converge at an ordinary plus
point by continuity of $N$, cannot converge on any joined component
because every original boundary collar has no interior zero, and
cannot converge at a point in the open minus copy or its compactified
end. Such a sequence is impossible in the compact Euclidean
complement. Thus $\mathcal E$ has no interior boundary in $\Omega$.

It is a nonempty open and closed subset of connected $\Omega$, and
hence equals $\Omega$. Vacuum now follows from
Proposition~\ref{n3:variation:adjoint}, since $N>0$ everywhere in
$\Omega$.
\end{proof}

\begin{lemma}[Equality forces a single spherical boundary component]
\label{n3:static:one-boundary-component}
In the Euclidean conformal double, every original component $S_i$
is a round sphere of radius $a=(8\kappa)^{-1}$ and has original area
$|S_i|_g=16\pi m^2=|S|_g$. Consequently $\ell=1$.
\end{lemma}

\begin{proof}
The completed conformal double in the proof of
Lemma~\ref{n3:static:no-internal-null} is Euclidean. In particular its
plus metric $q=h_+$ is flat up to every original boundary component.
Put $\phi=(1+\lambda)/2$. On each $S_i$,
Lemma~\ref{n3:static:base-completion} gives
\[
\phi=\tfrac12,\qquad
\partial_{\nu_h}\log\phi=\kappa,\qquad
\mathrm{II}_h=0,\qquad h|_{TS_i}=g|_{TS_i}.
\]
The conformal second-fundamental-form formula, with normal pointing
into the plus side, is therefore
\[
\mathrm{II}_q
 =\phi^2\bigl(\mathrm{II}_h+
 2\partial_{\nu_h}\log\phi\,h|_{TS_i}\bigr)
 =8\kappa\,q|_{TS_i},
 \qquad q|_{TS_i}=\tfrac1{16}g|_{TS_i}.
\]
Let $F_i$ be the Euclidean position vector on $S_i$ and $\nu_i$
its unit normal into the plus side. Its shape operator is
$8\kappa I$, so the tangential differential of $F_i-a\nu_i$ is
zero for $a=(8\kappa)^{-1}$. Connectedness makes $F_i-a\nu_i$
a constant point $c_i$. Thus the image lies on the sphere of radius
$a$ centered at $c_i$, with $\nu_i$ its outward radial normal.
The image is open in that sphere by local invertibility and closed
by compactness, hence is the whole sphere. The Euclidean embedding
therefore identifies $S_i$ diffeomorphically with it.
Since $q|_{TS_i}=g|_{TS_i}/16$, its area satisfies
\begin{equation}
\label{n3:static:individual-boundary-area}
|S_i|_g=16|S_i|_q=64\pi a^2
 =\frac{\pi}{\kappa^2}=16\pi m^2=|S|_g.
\end{equation}
Summing over the finite family yields
$|S|_g=\sum_{i=1}^{\ell}|S_i|_g=\ell|S|_g$. The boundary is nonempty
and $|S|_g>0$, hence $\ell=1$. In particular no connectedness or spherical
topology assumption was needed before this step.
\end{proof}

\begin{lemma}[Identification of the base and its mass]
\label{n3:static:schwarzschild-base}
The regular completion of \((\Omega,h,\lambda)\) is the
Schwarzschild static exterior
\begin{equation}
\label{n3:static:isotropic-base}
h=\left(1+\frac a\rho\right)^4\delta,\qquad
\lambda=\frac{1-a/\rho}{1+a/\rho},\qquad
\rho\geq a,
\end{equation}
where \(a=m/2\).
The completion is homeomorphic to the original \(\Mext\),
and \(S\) is identified diffeomorphically with its horizon sphere
in the regular base boundary structure.
\end{lemma}

\begin{proof}
Lemma~\ref{n3:static:one-boundary-component} identifies the boundary
in the Euclidean conformal double as one round sphere of radius
$a=(8\kappa)^{-1}$. The plus side is its outside: its normal is the
outward radial normal, its entire boundary is this sphere, and it
contains the distinguished unbounded end. Center Euclidean
coordinates at this sphere and write $\rho$ for their radius.

In the resulting Euclidean coordinates, write
\(\psi=2/(1+\lambda)\), so \(h=\psi^4\delta\).
The scalar-flat conformal equation gives \(\Delta_\delta\psi=0\)
outside the sphere, with \(\psi=2\) on it and \(\psi\to1\)
at infinity. The function \(1+a/\rho\) has the same data.
Their difference vanishes by the exterior maximum principle.
This proves Equation~\eqref{n3:static:isotropic-base}.
Its Schwarzschild mass is $2a=(4\kappa)^{-1}=m$, since
$\kappa=1/(4m)$. In particular $a=m/2$.

The regular base collar has the same underlying boundary topology
as the original collar, even when its coordinate is
\(\lambda=\sqrt N\) instead of the original defining function \(N\).
Together with the unchanged interior this identifies its completion
homeomorphically with \(\Mext\).
An interior base isometry extends uniquely to the metric completions.
On the regular boundary structures it is smooth: its boundary
restriction preserves induced lengths, hence is a smooth boundary
isometry, and normal geodesic coordinates extend it smoothly to
a collar.
Thus the completion has the claimed global topology, and
\(\Omega\cong(a,\infty)\times S^2\) is simply connected.
\end{proof}

\subsection{The time graph and the Kruskal attachment}
\label{n3:static:embedding-subsection}

\begin{proof}[Proof of Proposition~\ref{n3:static:classification}]
Lemma~\ref{n3:static:one-boundary-component} proves $\ell=1$ and the
spherical topology of $S$. Lemma~\ref{n3:static:no-internal-null} proves
\(N>0\) on \(\Omega\), and
Lemma~\ref{n3:static:schwarzschild-base} gives the global static
base with mass \(m\).
For the reconstruction write \(M=m\), so \(\kappa=1/(4M)\).
Simple connectivity and Lemma~\ref{n3:static:vanishing-twist}
give a global smooth function \(H\) with
\begin{equation}
\label{n3:static:time-primitive}
\dd H=-A.
\end{equation}
In static Schwarzschild coordinates use time \(T=t+H(x)\).
Equations~\eqref{n3:static:base-definitions} and
\eqref{n3:static:time-primitive} give the exact identity
\begin{equation}
\label{n3:static:reconstructed-metric}
-N\dd T^2+h
=-u^2\dd t^2+g_{ij}(\dd x^i+X^i\dd t)
                         (\dd x^j+X^j\dd t).
\end{equation}
The graph \(T=H(x)\), corresponding to \(t=0\), therefore induces
the original metric \(g\).
The stationary future unit normal is \(n=(\partial_t-X)/u\).
With the second-fundamental-form convention of the theorem,
\begin{equation}
\label{n3:static:reconstructed-k}
K_{\mathrm{graph}}
=\frac1{2u}\bigl(\partial_tg-\mathcal L_Xg\bigr)
=-\frac1{2u}\mathcal L_Xg
=K.
\end{equation}
This calculation recovers the original \(K\), including a
non-time-symmetric slice.

We next establish smoothness in the original boundary structure.
If \(u|_S=0\), Equation~\eqref{n3:static:double-norm} gives
\[
A=\frac{Y^\flat}{a^2-s^2|Y|_g^2}.
\]
This is smooth in the original collar.
The primitive \(H\) therefore extends smoothly and finitely to \(S\).
For instance, fix its values on one interior collar section and
integrate its smooth normal derivative to the boundary.
Tangential derivatives of this extension agree with those prescribed
by \(\dd H=-A\), by continuity from the interior.

If \(u|_S>0\), the numerator
\(X^\flat-(2\kappa)^{-1}\dd N\) vanishes as a one-form on \(S\),
by Equation~\eqref{n3:static:simple-norm}.
Smooth division by the defining function \(N\) gives
\begin{equation}
\label{n3:static:log-time}
A=\frac1{2\kappa}\dd\log N+\beta,\qquad
H=-\frac1{2\kappa}\log N+B,
\end{equation}
where \(\beta\) and \(B\) are smooth in the original collar.
The smoothness of \(B\) follows by applying the same normal-integration
argument to its differential \(-\beta\).
In this case \(H\to+\infty\) at the boundary, so static time itself
is not a regular attachment coordinate.

Let \(r\) be the Schwarzschild area radius. On the whole exterior
\[
N=1-\frac{2M}{r},\qquad
r=\frac{2M}{1-N}.
\]
Normalize the tortoise coordinate and Kruskal coordinates by
\begin{equation}
\label{n3:static:kruskal-definitions}
\begin{aligned}
r_*&=r+2M\log\left(\frac r{2M}-1\right),\\
\mathsf U&=-e^{-\kappa(T-r_*)},&
\mathsf V&= e^{\kappa(T+r_*)}.
\end{aligned}
\end{equation}
On the right exterior \(\mathsf U<0,\mathsf V>0\), and
the Schwarzschild metric extends as
\[
\mathbf g_M
=-\frac{32M^3}{r}e^{-r/(2M)}
   \dd\mathsf U\,\dd\mathsf V+r^2\dd\omega^2.
\]
An exact useful form of Equation~\eqref{n3:static:kruskal-definitions}
is
\begin{equation}
\label{n3:static:kruskal-factor}
\mathsf U=-Q(N)\sqrt N\,e^{-\kappa T},\qquad
\mathsf V= Q(N)\sqrt N\,e^{\kappa T},\qquad
Q(N)=\frac{\exp\!\bigl(1/[2(1-N)]\bigr)}{\sqrt{1-N}}.
\end{equation}
The function \(Q\) is smooth and positive near zero, with
\(Q(0)=\sqrt e\).

If \(u|_S=0\), both \(H\) and
\(\sqrt N=s\sqrt{a^2-s^2|Y|^2}\) are smooth.
Substitution into Equation~\eqref{n3:static:kruskal-factor} gives
\(\mathsf U=\mathsf V=0\) on \(S\) and
\[
\partial_s\mathsf U|_S
=-\sqrt e\,\kappa e^{-\kappa H|_S}<0,\qquad
\partial_s\mathsf V|_S
=\sqrt e\,\kappa e^{\kappa H|_S}>0.
\]
This is a smooth attachment at the bifurcation sphere.
If \(u|_S>0\), Equation~\eqref{n3:static:log-time} instead gives
\begin{equation}
\label{n3:static:future-attachment}
\mathsf U=-Q(N)N e^{-\kappa B},\qquad
\mathsf V=Q(N)e^{\kappa B}.
\end{equation}
These are smooth in the original collar;
\(\mathsf U\) vanishes simply and
\(\mathsf V|_S=\sqrt e\,e^{\kappa B|_S}>0\).
Thus the attachment is on the future horizon, with each generator
met once after the angular identification.

The angular regularity requires a further check when
\(u|_S>0\), because the regular base coordinate was \(\lambda\),
not the original \(N\).
On the smooth reflected collar in
Equation~\eqref{n3:static:even-base}, normal projection onto its fixed
boundary is invariant under \(\lambda\mapsto-\lambda\).
This follows from uniqueness of the normal geodesics and reflection
invariance of the metric.
Schwarzschild angular coordinates are constant along these normal
geodesics, and their boundary values are the smooth round-sphere
identification from Lemma~\ref{n3:static:schwarzschild-base}.
They are therefore smooth even functions of \(\lambda\).
A smooth even function, with smooth tangential parameters, is a
smooth function of \(N=\lambda^2\) for \(N\geq0\):
all odd Taylor coefficients vanish, and repeated application of
\((2\lambda)^{-1}\partial_\lambda\) to the Taylor remainder gives
continuous derivatives of every order at zero.
This proves the required angular smoothness in the original collar.
In the zero boundary-lapse case the original collar is already a
smooth base collar, so its completed base isometry supplies the
angular extension directly.

We have constructed a smooth map \(\iota:\Mext\to\) Schwarzschild.
Equation~\eqref{n3:static:reconstructed-metric} extends to \(S\)
by continuity, hence \(\iota^*\mathbf g_M=g\) there as well.
The positive definiteness of \(g\) proves that \(\dd\iota\) is
injective at every boundary point.
Interior injectivity follows from the global base isometry;
boundary injectivity follows from the sphere identification.
The boundary image, at \(r=2M\), is disjoint from the open image,
where \(r>2M\).
The map is proper: a compact subset of Schwarzschild has bounded
area radius, and its inverse image is closed in the compact base
region corresponding to
\([2M,R]\times S^2\).
Thus the injective immersion is an embedding.

The future normal extends smoothly even where \(u=0\).
To avoid taking a limit of the expression involving \(u^{-1}\),
choose the smooth future timelike Kruskal vector
\(Z=\partial_{\mathsf U}+\partial_{\mathsf V}\) near the boundary.
Let \(Z^\top\) be its tangential projection along the embedded
spacelike hypersurface and set
\[
\widehat n=
\frac{Z-Z^\top}
{\sqrt{-\mathbf g_M(Z-Z^\top,Z-Z^\top)}}.
\]
The denominator is positive, and this is a smooth future unit
normal. On the open exterior it agrees with the normal used in
Equation~\eqref{n3:static:reconstructed-k}.
Its second fundamental form is smooth and equals \(K\) in the
interior, so equality extends to \(S\).

There is also a smooth spacelike extension of the image through
its boundary. In the positive boundary-lapse case, use
\((\mathsf U,\omega)\) as hypersurface coordinates; its normal
derivative is nonzero by
Equation~\eqref{n3:static:future-attachment}.
The remaining null coordinate is a smooth graph function and can
be extended across \(\mathsf U=0\).
In the zero boundary-lapse case use
\((\chi,\omega)\), where
\(\chi=(\mathsf V-\mathsf U)/2\); the displayed normal derivatives
give \(\partial_s\chi>0\).
Again extend the remaining graph function across zero.
Smooth one-sided graph functions extend across a smooth boundary,
and compactness of \(S\), together with openness of the spacelike
condition, preserves spacelikeness on a sufficiently short collar.
This extension makes no assertion about the original data behind
\(S\).

Finally, the exact inverse-metric identity gives
\begin{equation}
\label{n3:static:end-slope}
|\dd H|_h^2
=\left|\frac{X^\flat}{N}\right|_C^2
=\frac{|X|_g^2}{u^2N}
\longrightarrow\frac{|b|^2}{(b^0)^2}<1.
\end{equation}
Choose a constant \(a_0<1\) larger than the limiting slope norm.
Along radial rays in the Schwarzschild base, uniformly in angle,
\[
|H(r,\omega)|
\leq C+a_0\int_R^r
\frac{\dd\rho}{\sqrt{1-2M/\rho}}
=a_0r+O(\log r).
\]
Since \(r_*=r+O(\log r)\), it follows that
\[
H-r_*\longrightarrow-\infty,\qquad
H+r_*\longrightarrow+\infty.
\]
This is approach to the chosen Schwarzschild spatial infinity.
It allows a nonzero asymptotic tilt and does not impose \(P=0\)
or \(K=0\).
All conclusions of Proposition~\ref{n3:static:classification} follow.
\end{proof}

Propositions~\ref{n3:variation:adjoint} and~\ref{n3:static:classification}
prove Theorem~\ref{n3:intro:rigidity}.
